\setcounter{tocdepth}{2}

%%%%%%%%%%%%%%%%%%%%%%%%%%%%%%%%%%%% TITULO TOC %%%%%%%%%%%%%%%%%%%%%%%%%%%%%%%%%%%%%%%%

\setlength{\cftbeforetoctitleskip}{0em}
% \setlength{\cftaftertoctitleskip}{7em}
\newcommand{\tocdisplaytitle}[2][1.2cm]{%
    \hfill\fontfamily{lmr}\selectfont\resizebox{!}{#1}{#2}%
}
\renewcommand{\cfttoctitlefont}{\tocdisplaytitle}
\renewcommand{\cftaftertoctitle}{
    \resizebox{!}{1.2cm}{\,\colorbox{rojoCIMAT}{\!\vphantom{I}\!}}
}

% Las entradas sin número no reciben el espacio que normalmente ocupa el número.
% Se agrega un margen breve dentro del recuadro gris y se conserva el texto del marcador.
\newcommand{\chaptertocunnumbered}[1]{%
    \texorpdfstring{\hspace*{0.5em}#1}{#1}%
}

%%%%%%%%%%%%%%%%%%%%%%%%%%%%%%%%%%%% SEPARACIÓN %%%%%%%%%%%%%%%%%%%%%%%%%%%%%%%%%%%%%%%%

% Capitulos
\setlength{\cftchapnumwidth}{2.5em}
\setlength{\cftbeforechapskip}{0.7em} % arriba
\renewcommand{\cftchapafterpnum}{\chaptertocboxend\vspace{0.5em}} % abajo
% Secciones
\setlength{\cftsecnumwidth}{\cftchapnumwidth}   
\setlength{\cftsecindent}{1.5em}
\setlength{\cftbeforesecskip}{0.2em} % arriba
\renewcommand{\cftsecafterpnum}{\vspace{0.1em}} % abajo
% Subsecciones
\setlength{\cftsubsecnumwidth}{4em} % hueco para tres números «5.14.1»
\setlength{\cftsubsecindent}{\dimexpr\cftsecindent+\cftsecnumwidth\relax} % indentado
\setlength{\cftbeforesubsecskip}{0.1em}
% Subsubsecciones
\setlength{\cftsubsubsecnumwidth}{3.6em}
\setlength{\cftsubsubsecindent}{\dimexpr\cftsubsecindent+\cftsubsecnumwidth\relax}

%%%%%%%%%%%%%%%%%%%%%%%%%%%%%%%%%%%%% CAPÍTULOS %%%%%%%%%%%%%%%%%%%%%%%%%%%%%%%%%%%%%%%%

\cftpagenumbersoff{chap}

% Contorno de la entrada: las dos marcas permiten títulos de varias líneas
% conservando el margen, la tipografía y la altura original en una sola línea.
\newcounter{chaptertocentry}
\newcommand{\chaptertocbox}{%
    \makebox[0pt][l]{%
        \tikz[remember picture,overlay]{%
            \coordinate (toc-start-\chaptertocid) at (0,0);
        }%
    }%
}
\newcommand{\chaptertocboxend}{%
    \begingroup\large
    \tikz[remember picture,overlay]{%
        \coordinate (toc-end-\chaptertocid) at (0,0);
        \coordinate (toc-bottom-\chaptertocid) at
            (toc-start-\chaptertocid |- toc-end-\chaptertocid);
        \draw[color=grisCIMAT, line width=1pt]
            ([yshift=0.95em]toc-start-\chaptertocid)
            rectangle
            ([xshift=\textwidth,yshift=-0.27em]toc-bottom-\chaptertocid);
    }%
    \endgroup
}

% Recuadro del número, de ancho automático
\newcommand{\chaptertocnumberline}[1]{%
    \makebox[\cftchapnumwidth][l]{%
        \begingroup
        \setlength{\fboxsep}{0pt}%
        \colorbox{grisCIMAT}{%
            \hspace{0.15em}%
            \rule[-0.27em]{0pt}{1.22em}%
            \bfseries\color{white}#1\textordfeminine
            \hspace{0.15em}%
        }%
        \endgroup
    }%
    {Sesión:\enspace}%
}

% Formato de las entradas
\renewcommand{\cftchapfont}{\color{grisCIMAT!75!black}\large\chaptertocbox}
% \renewcommand{\cftchapaftersnumb}{%
%     {\color{grisCIMAT!80!black}Clase:\enspace}%
% }

% Aplicamos el formato únicamente a los capítulos
\makeatletter
\let\oldl@chapter\l@chapter
\renewcommand{\l@chapter}[2]{%
    \begingroup
    \stepcounter{chaptertocentry}%
    \edef\chaptertocid{\arabic{chaptertocentry}}%
    \let\numberline\chaptertocnumberline
    \oldl@chapter{#1}{#2}%
    \endgroup
}
\makeatother

%%%%%%%%%%%%%%%%%%%%%%%%%%%%%%%%%%%%% SECCIONES %%%%%%%%%%%%%%%%%%%%%%%%%%%%%%%%%%%%%%%%

\renewcommand{\cftsecfont}{\fontfamily{bch}\selectfont\color{grisCIMAT!0!black}}
\renewcommand{\cftsecaftersnum}{.}

\renewcommand{\cftsubsecfont}{\fontfamily{bch}\selectfont\color{grisCIMAT!0!black}}
\renewcommand{\cftsubsecaftersnum}{.}

% \renewcommand{\cftsubsubsecfont}{\small\fontfamily{cmss}\selectfont}
